\chapter*{Итог: из чего складывается переориентация}
\addcontentsline{toc}{chapter}{Итог: из чего складывается переориентация}
\markboth{Итог}{Итог}

Соберём главы вместе.
\begin{enumerate}
\item \textbf{Из чего сделан материал} --- ГПУ цирконий, водород в тетраэдрических пустотах, $\delta$-гидрид с ГЦК
  подрешёткой, которая получается из ГПУ сдвигами частичных дислокаций (главы~\ref{ch:atoms}--\ref{ch:defects}).
\item \textbf{Как ориентированы зёрна} --- текстура трубы: оси $c$ у радиуса с наклоном к окружности, факторы Кернса
  и полное распределение ориентаций (глава~\ref{ch:texture}).
\item \textbf{Сколько выпадет и когда} --- линии растворимости TSSD и TSSP, гистерезис между ними, угол диаграммы
  Zr--H (главы~\ref{ch:gibbs}, \ref{ch:solutions}, \ref{ch:diagrams}).
\item \textbf{Где окажется водород} --- быстрая диффузия выравнивает концентрацию на масштабе зёрен. Водород
  стекается к растянутым местам: торцам пластинок, вершинам трещин (глава~\ref{ch:diffusion}).
\item \textbf{Какой формы и в какой плоскости} --- тонкие пластинки около базиса: так меньше упругая энергия
  несоответствия (главы~\ref{ch:lattice}, \ref{ch:elastic}).
\item \textbf{В каких зёрнах и когда зародится} --- барьер зарождения, в показатель которого входят упругая
  энергия, работа нагрузки, дефекты и границы (главы~\ref{ch:elastic}, \ref{ch:nucleation}).
\item \textbf{Как пластинки выстраиваются} --- поле соседей, пластическая аккомодация, память о прежних гидридах
  (главы~\ref{ch:elastic}, \ref{ch:plasticity}).
\item \textbf{Какая доля станет радиальной} --- конкуренция зёрен с разной ориентацией, то есть текстура
  (глава~\ref{ch:texture}).
\item \textbf{С чем сравнивать} --- травленые шлифы, ДСК, дифракция, EBSD: что именно они меряют и где искажают
  картину (глава~\ref{ch:methods}).
\item \textbf{Чем это грозит} --- радиальные гидриды как готовые трещины, переход хрупкость--пластичность,
  замедленное растрескивание (глава~\ref{ch:mech}).
\end{enumerate}
Модель гидридов --- это численная реализация пунктов 3--8 одновременно, проверяемая по пункту 9 ради ответа на
пункт 10. Читать её код и учебник по ней (\texttt{../textbook}) после этого фундамента должно быть заметно легче.
